\chapter{Literature Review}
\label{ch:literature}

This chapter reviews the foundational concepts and related work underpinning this research.

\section{First-Order Logic and Deontic Extensions}

\subsection{First-Order Logic (FOL)}

First-Order Logic provides a formal language for representing knowledge with quantifiers over individuals \cite{brachman2004knowledge}. The key components include:

\begin{itemize}
    \item \textbf{Constants}: Named individuals (e.g., \texttt{John}, \texttt{AIT})
    \item \textbf{Predicates}: Properties and relationships (e.g., \texttt{Student(x)}, \texttt{Enrolled(x,y)})
    \item \textbf{Quantifiers}: Universal ($\forall$) and existential ($\exists$)
    \item \textbf{Connectives}: Conjunction ($\wedge$), disjunction ($\vee$), implication ($\rightarrow$), negation ($\neg$)
\end{itemize}

FOL is decidable for certain fragments and provides a balance between expressiveness and computability \cite{russell2016artificial}.

\subsection{Deontic Logic}

Deontic logic extends classical logic with normative operators \cite{governatori2010defeasible}:

\begin{itemize}
    \item \textbf{O($\phi$)}: Obligation --- $\phi$ must be the case
    \item \textbf{P($\phi$)}: Permission --- $\phi$ may be the case
    \item \textbf{F($\phi$)}: Prohibition --- $\phi$ is forbidden (equivalent to O($\neg\phi$))
\end{itemize}

Standard Deontic Logic (SDL) satisfies the axiom $O(\phi) \rightarrow P(\phi)$, meaning all obligations imply permissions \cite{mcnamara2006deontic}.

\subsection{Higher-Order Logic Considerations}

Second-Order Logic (SOL) extends FOL by allowing quantification over predicates:
$$\forall P. P(x) \rightarrow \ldots$$

Higher-Order Logic (HOL) treats functions as first-class objects with lambda expressions \cite{farmer2008seven}. While more expressive, both introduce undecidability concerns. As Brachman and Levesque note, ``the price of greater expressiveness is often computational intractability'' \cite{brachman2004knowledge}.

\section{Large Language Models for Text Classification}

\subsection{Transformer Architecture}

Modern LLMs are based on the Transformer architecture \cite{vaswani2017attention}, which uses self-attention mechanisms to process sequential data. Key innovations include:

\begin{itemize}
    \item Multi-head attention for parallel processing
    \item Positional encodings for sequence order
    \item Layer normalization for training stability
\end{itemize}

\subsection{Instruction-Tuned Models}

Recent work has shown that instruction-tuned models outperform base models on classification tasks:

\begin{itemize}
    \item \textbf{Mistral 7B} \cite{jiang2023mistral}: Achieves competitive performance with 7B parameters through grouped-query attention and sliding window attention
    \item \textbf{Llama 3} \cite{touvron2023llama}: Meta's open-source family with strong text classification capabilities
    \item \textbf{Phi-3} \cite{abdin2024phi}: Microsoft's compact model demonstrating that careful data curation can compensate for smaller size
\end{itemize}

\subsection{Classification Performance}

Studies have demonstrated LLM effectiveness for legal and policy text classification:

\begin{quote}
``Fine-tuned Llama 3 outperforms RoBERTa-large on text classification tasks while maintaining generalization'' \cite{meta2024llama}
\end{quote}

\section{SHACL and Semantic Web Technologies}

\subsection{RDF and Knowledge Graphs}

The Resource Description Framework (RDF) provides a graph-based data model with subject-predicate-object triples \cite{w3c2014rdf}:

\begin{verbatim}
:Student123 rdf:type :Student .
:Student123 :enrolled :Course456 .
\end{verbatim}

\subsection{SHACL Constraints}

SHACL (Shapes Constraint Language) defines validation constraints over RDF graphs \cite{w3c2017shacl}:

\begin{itemize}
    \item \textbf{Node Shapes}: Target specific node types
    \item \textbf{Property Shapes}: Constrain property values
    \item \textbf{Severity Levels}: Violation, Warning, Info
\end{itemize}

Example SHACL shape:
\begin{verbatim}
ex:StudentShape a sh:NodeShape ;
    sh:targetClass ex:Student ;
    sh:property [
        sh:path ex:paid ;
        sh:minCount 1 ;
    ] .
\end{verbatim}

\subsection{SHACL for Policy Compliance}

Palmirani and Governatori demonstrate the adequacy of SHACL for policy formalization:

\begin{quote}
``SHACL's constraint language, based on first-order patterns, provides adequate expressiveness for policy rules in institutional contexts'' \cite{palmirani2018from}
\end{quote}

\section{Related Work}

\subsection{Automated Policy Extraction}

Prior work on policy extraction includes:

\begin{itemize}
    \item \textbf{Rule-based approaches}: Pattern matching for deontic markers \cite{maxwell2012textual}
    \item \textbf{Machine learning}: SVM and neural classifiers for policy identification \cite{brodie2006automated}
    \item \textbf{Hybrid methods}: Combining NLP with ontological reasoning \cite{wyner2012rule}
\end{itemize}

\subsection{Logic Formalization of Policies}

Approaches to formal policy representation:

\begin{itemize}
    \item \textbf{Defeasible Logic}: Handling exceptions and conflicts \cite{governatori2010defeasible}
    \item \textbf{Event Calculus}: Temporal reasoning about actions \cite{kowalski1989logic}
    \item \textbf{Description Logic}: Ontology-based constraints \cite{baader2003description}
\end{itemize}

\subsection{LLMs for Legal NLP}

Recent applications of LLMs to legal text:

\begin{itemize}
    \item Contract analysis and clause extraction \cite{hendrycks2021cuad}
    \item Legal question answering \cite{zhong2020nlpcc}
    \item Statute understanding \cite{niklaus2021lextreme}
\end{itemize}

\section{Gap Analysis}

While prior work addresses individual components, there is limited research combining:

\begin{enumerate}
    \item LLM-based policy rule identification
    \item Deontic logic formalization
    \item SHACL-based validation
\end{enumerate}

This thesis addresses this gap by providing an end-to-end pipeline from natural language policies to executable validation constraints.
